\chapter{Introduction}
\label{ch:introduction}

\section{Background and Motivation}

Large language models (LLMs) have become strong general-purpose generators, but the knowledge encoded in their parameters is necessarily incomplete and fixed at training time. When a question concerns recent events, long-tail entities, or information outside the training distribution, an LLM may still produce a fluent answer even when the underlying facts are unsupported. This behaviour limits the reliability of LLMs in knowledge-intensive applications~\cite{hallucination_survey}.

Retrieval-Augmented Generation (RAG) mitigates this problem by retrieving external passages and conditioning generation on the resulting evidence~\cite{rag_survey}. Because the evidence source can be updated independently of the language model, RAG provides a practical way to extend the model's accessible knowledge without repeated model training. A fixed retrieval strategy, however, is not equally suitable for every question. For a simple question that the model can already answer, irrelevant passages can introduce distractors and additional latency. For a complex multi-hop question, one retrieval query based on the original question may fail to expose the intermediate entity needed to locate the next supporting passage~\cite{ren2023,ircot}.

Adaptive RAG addresses this mismatch by changing the retrieval strategy according to the question or the model's state. Existing approaches include question-level complexity and self-knowledge routers~\cite{adaptiverag,skr}, uncertainty-triggered retrieval during generation~\cite{flare,dragin}, and learned policies that decide retrieval actions at individual reasoning steps~\cite{selfrag,deeprag}. These approaches demonstrate that retrieval need is query-dependent, but each has a different limitation. A question-level route is decided before the intermediate reasoning states are known; uncertainty is an imperfect proxy for correctness; and a learned step-level policy can require additional trajectory construction, supervision, or preference optimisation.

Multi-hop question answering introduces a second decision that is closely coupled with retrieval: whether to decompose the question. Question decomposition can convert a compositional question into shorter sub-questions that expose bridge entities and form more precise retrieval queries~\cite{decomprc,break_qdmr}. Nevertheless, decomposition is not uniformly reliable. An incorrect intermediate answer can misdirect later queries, while a short sub-question can omit comparison targets, entity pairs, or disambiguating conditions from the original question. Forcing every question through a decomposed path may therefore discard useful global context even when the retrieved passages themselves are relevant.

This thesis consequently frames adaptive RAG as a \emph{decomposition-routing problem}. A system should decide at two levels: whether each decomposed sub-question requires external evidence, and whether the final answer should be taken from the original-question path or the decomposed reasoning path. Preserving both paths makes decomposition a candidate strategy rather than an irreversible early commitment. This perspective motivates SEEK (Self-certainty Evaluated Evidence-gated Knowledge-retrieval), a lightweight framework that trains a compact decomposer but keeps retrieval gating and final path verification transparent and training-free.

Accordingly, this thesis investigates whether competitive adaptive RAG can be achieved with lightweight decomposition training and interpretable inference-time routing, without learning a separate retrieval policy. Its objectives are to produce structurally reliable and executable sub-questions, examine both the value and the limits of token-level self-certainty for sub-question retrieval, design a context-preserving dual-path procedure with transparent final verification, and evaluate the resulting framework across single-hop and multi-hop open-domain question answering in terms of accuracy, retrieval behaviour, offline adaptation cost, and failure modes.

\section{Research Contributions}

Our contributions are as follows:

\begin{enumerate}
    \item \textbf{Dual-path decomposition routing.} This thesis treats adaptive RAG as a choice between an original-question answer path and a decomposed answer path, rather than only a retrieve-or-skip decision. SEEK keeps both candidates until a final verification step, so decomposition becomes a reversible strategy instead of an early commitment. The two paths prove complementary: in the final traces, an oracle choice between them would add a further 2.9 to 6.0 EM, and the training-free verifier recovers part of this headroom.

    \item \textbf{A low-cost executable decomposer.} The decomposer is adapted in a single supervised LoRA~\cite{lora} pass on 12{,}000 examples and emits KG-style \texttt{entity >> relation} sub-questions with explicit \texttt{\#N} dependencies that resolve into directly executable retrieval queries. Keeping original-question evidence available at every decomposed step preserves global context that short sub-questions may omit, without changing the decomposer or training another retriever.

    \item \textbf{A simple confidence gate with a measured accuracy--retrieval trade-off.} Tail Leakage, computed from the probability mass the model keeps within its own top-five token candidates during a single greedy pass, gates sub-question retrieval. It captures a simple pattern: answers are more likely to be wrong when more probability falls outside these leading candidates. The score therefore measures only the total leaked mass, regardless of how probability is divided within the top five or distributed across the tail.

    \item \textbf{Training-free late verification.} A deterministic verifier reads direct-answer confidence, decomposition stability, and answer-type compatibility, and selects the final path so that every routing decision remains inspectable. The oracle comparison quantifies the remaining selection headroom without attributing candidate-generation gains to the verifier itself.
\end{enumerate}
